% Chapter 1 -- Introduction
% Standalone, compilable draft (FPT-template article class + booktabs), sibling to
% sec_2_background.tex / sec_8_discussion.tex. Frames the problem, the modular
% approach, the three research questions, the claimed contributions, the scope, and
% the thesis outline. Its load-bearing job is to OWN the proposal drift: three
% promises in docs/report/thesis_proposal.tex were narrowed or dropped during the
% research, and this chapter states each as a deliberate, evidence-backed decision
% rather than letting a committee surface it first.
%
% MERGE/NUMBERING NOTE: this file opens with \section{Introduction} for standalone
% compilation and numbering, but at final assembly the assembler promotes it to a
% chapter (Chapter 1). Forward references to sibling chapters use typed numerals
% ("Chapter~3", "Section~7.4") because the chapter files do not share a label
% namespace when built separately. Internal experiment records are cited as plain
% text "Finding~\#NN" (entries in docs/findings.md), the convention used by the
% method/results/discussion chapters; published literature uses BibTeX \cite{}.
%
% CONSISTENCY CONSTRAINTS (checked at draft time):
%   - Every claim here is a subset of the THESIS_PLAN Section-6 claims map (C1--C11).
%     Nothing is stated beyond the PARTIALLY-HOLDS scope of the held-out completion
%     result (C9). The three claimed contributions match sec_2_background.tex Sec 2.5
%     exactly (donor metric / recall root-cause / synthetic-pretraining redundancy);
%     the pipeline, runtime characterisation, and amodal-GT builder are named as
%     ENGINEERING contributions, not scientific claims, and tracking is explicitly
%     not a contribution (as Sec 2.5 states outright).
%   - Proposal-drift citations: multi-class->car-only cites Finding~\#20 (classifier
%     only); CD+EMD->CD-only cites the domain-gap bundle (#15--19 + the input-
%     distribution diagnosis + the KITTI-like fix #26), NOT #20; two-stage-training
%     "reduces annotation" is reframed as the #25/#30 negative finding.
%   - "offline" is a SCOPE statement, not a walk-back: the proposal is silent on
%     runtime, so it is not presented as a dropped promise.
%
% Honesty-phrasing rules applied (personal_agenda P1):
%   - recall is stated against its car-only, >=10-surviving-points denominator
%   - the held-out completion result is PARTIALLY HOLDS, never unqualified
%   - negative results are named as findings, not buried
%   - the detection headline is a single-sequence (seq 08) result, stated as such
%
\documentclass[a4paper,12pt]{article}
\usepackage[english]{babel}
\usepackage[utf8]{inputenc}
\usepackage[T1]{fontenc}
\usepackage{lmodern}
\usepackage{amsmath,amssymb}
\usepackage{booktabs}
\usepackage{float}
\usepackage{siunitx}
\usepackage{microtype}
\usepackage[margin=2.5cm]{geometry}
\usepackage{graphicx}
\usepackage{multirow}
\usepackage{url}

\begin{document}

\section{Introduction}
\label{ch:introduction}

\subsection{Problem and motivation}
\label{sec:intro-motivation}

Perception for autonomous driving rests on a LiDAR scanner that returns, several
times a second, an unordered cloud of over a hundred thousand three-dimensional points.
Turning that raw geometry into a list of objects --- where the cars are, how large
they are, which way they face --- is the first step in almost every downstream
driving task. The dominant way to do this is to train a deep network end to end on
densely annotated scans. The public benchmarks that make such training possible make
the cost of the approach plain: SemanticKITTI~\cite{Behley_2019_ICCV}, built on the
KITTI odometry data~\cite{geiger2012we}, provides a dense point-wise semantic label
for every point of its labelled sequences (billions of points in total), an
annotation effort that is far out of reach for most groups and datasets. A model trained this way is also
difficult to inspect: when it misses a car, the failure is distributed across
millions of weights, and there is no single stage to point at.

This thesis takes the opposite starting point. It asks how far a \emph{modular},
largely annotation-free pipeline can go using classical geometry for the parts that
geometry solves well and small learned models only where discrimination genuinely
needs them. Ground points are removed and the remaining cloud is clustered into
candidate objects by unsupervised density clustering; a lightweight binary
classifier decides which clusters are cars; a short tracker filters transient
detections; and, for each detected car, a point-completion network fills in the
occluded surface that a single viewpoint never sees. Every stage is separately
observable and separately debuggable, and the only learned components --- the
classifier and the completion network --- are trained without any point-wise
semantic labels. The price of this transparency, and the limits it runs into, are as
much the subject of this thesis as its successes.

A second thread runs alongside detection. Completing a partially observed car ---
recovering the far side that the sensor cannot reach --- is attractive for
downstream box estimation, and point-completion networks such as
PCN~\cite{yuan2019pcnpointcompletionnetwork} are well studied on synthetic
benchmarks. Applying them to real LiDAR raises a problem that turns out to be
central: there is no clean ground-truth complete shape to measure against, and the
obvious substitute --- accumulating the same static car across many frames to build
a pseudo-ground-truth --- is itself one-sided and, as this thesis shows, invalid as
a reference. Much of the completion half of the work is therefore about how to
\emph{measure} completion quality on real data at all, before asking whether it
helps.

\subsection{Approach}
\label{sec:intro-approach}

The system is a modular hybrid. A raw scan is filtered and downsampled, the ground
plane is removed by a relative-height RANSAC fit, and the remaining points are
grouped into candidate objects by HDBSCAN density clustering. A geometric filter
discards clusters that cannot be cars on size grounds, a dual-branch PointNet-style
binary classifier labels each surviving cluster as \emph{car} or \emph{not-car}, and
a centroid tracker links detections across frames so that transient, one-frame
clusters can be voted away. Each detected car is then passed to a PCN completion
network that predicts a filled point cloud for the whole vehicle from the single
partial view. The design commitment is deliberate: classical geometry does the
spatial decomposition it is reliable at, while the learned models are kept small and
are asked only to discriminate and to reconstruct. The full architecture, the
per-stage parameters, and the classifier's training procedure are the subject of
Chapter~3, whose pipeline diagram (Figure~3.1) shows the stages and their measured
per-stage cost; this chapter states only the shape of the approach.

Two points about the approach are worth making early, because they shape what the
thesis can and cannot claim. First, the tracker is used purely as a temporal filter
--- it removes flicker detections and does nothing to interpolate or recover missed
cars --- so it is not counted among the contributions. Second, the completion
network is trained only on synthetic partial renders made to imitate LiDAR
sampling; it never sees a labelled real car. This is what makes the pipeline
near-annotation-free, and it is also the source of the sim-to-real questions that
Chapters~4 and~7 take up.

\subsection{Research questions}
\label{sec:intro-rq}

The work is organised around three questions.

\begin{description}
  \item[RQ1.] How can the quality of point-cloud completion be evaluated on
    \emph{real} LiDAR, where no clean reference shape exists and an accumulated
    pseudo-ground-truth is itself one-sided? This is a measurement question, and it
    is prior to any claim that completion helps. It is answered in Chapter~7.
  \item[RQ2.] Does learned completion recover genuinely \emph{unseen} surface --- as
    opposed to redistributing the points it was already given --- and does it
    improve downstream amodal bounding boxes on real LiDAR cars? And does any such
    benefit survive on a held-out sequence under a pre-registered test? This is
    answered in Chapter~7, with the held-out outcome reported exactly as it fell.
  \item[RQ3.] What limits the recall of a training-free clustering-plus-classifier
    detector, and is synthetic pretraining actually necessary for the classifier?
    The first half is answered in Chapter~5, the second in Chapter~4.
\end{description}

\subsection{Contributions}
\label{sec:intro-contributions}

The thesis makes three claimed contributions, all methodological or diagnostic
rather than architectural. It does not claim novelty for the pipeline itself: the
clustering, classification, tracking, and completion stages are established building
blocks, reimplemented and tuned here.

\begin{enumerate}
  \item \textbf{A leakage-free metric for completion quality on real LiDAR.} The
    thesis first shows that reference-based Chamfer distance against an accumulated
    pseudo-ground-truth is invalid on real cars --- the raw partial input scores
    \emph{best} on that metric, because the reference is itself occluded on the same
    side (Finding~\#26). It then defines a donor-frame occluded-side coverage
    metric that scores how much of the surface seen only in \emph{other} frames of
    the same static car a completion recovers, with a hallucination guard against
    surface invented outside the car's box, and validates it against a four-item
    gate (Findings~\#32/\#37). Applied to real sequence-08 cars, the metric shows
    that PCN recovers occluded surface several times above the symmetry-mirror
    baseline and improves amodal boxes on static, gate-passed cars, and that the
    benefit partially holds on a pre-registered held-out sequence
    (Findings~\#29/\#42). This is the completion half's central contribution and is
    developed in Chapter~7.
  \item \textbf{A root-cause account of the recall bottleneck, with a chain of
    negative repair results.} Recall is limited not by the classifier but by
    density clustering splitting single cars into fragments (Finding~\#23). A
    search of repair strategies --- larger minimum cluster size, fragment merging,
    adaptive and BEV clustering, temporal accumulation --- fails to generalise,
    each for a structural reason (Findings~\#21/\#24); the one lever that helped
    was the clustering resolution itself (Finding~\#34). The negative results are reported as findings in their
    own right in Chapter~5.
  \item \textbf{Evidence that synthetic pretraining is redundant for the classifier
    at scale.} A from-scratch classifier matches a synthetically pretrained one once
    real training clusters are plentiful, and a cross-domain matrix shows the
    sim-to-real gap is total and symmetric (Findings~\#25/\#30). This directly
    replaces a contribution promised in the proposal, as discussed in
    Section~\ref{sec:intro-scope}, and is presented in Chapter~4.
\end{enumerate}

Alongside these, the thesis delivers several \emph{engineering} contributions that
are not claimed as scientific results: the modular pipeline and its runtime
characterisation, an amodal ground-truth box builder for the completion evaluation,
and the reproducibility infrastructure (deterministic evaluation, invariant tests)
that the frozen results rest on. A number of directions that did not work ---
architecture swaps and real-data fine-tuning for completion (Findings~\#16/\#17/\#19/\#28),
and a sparse-fallback length model (Finding~\#40) --- are reported as negative
findings rather than omitted.

\subsection{Scope and reconciliation with the proposal}
\label{sec:intro-scope}

The thesis is deliberately narrow, and the boundaries are worth stating plainly
rather than leaving to be inferred. \textbf{The detector is car-only:} the
classifier is binary (car / not-car), and no other class is detected. \textbf{The
system is offline:} it runs at about \SI{623.5}{\milli\second} per frame
(\num{1.6}~frames per second) on the reference hardware, and reducing this to
real-time is out of scope; the runtime is characterised (Section~4.5,
Findings~\#33/\#34), not optimised further. \textbf{The detection headline rests on
a single sequence:} sequence~08 (\num{4071} scans) is the reporting sequence, with
sequence~00 used for tuning and for a pre-registered held-out completion check;
per-sequence recall variation across the labelled set is therefore not broken out,
a breadth limitation stated openly in Chapter~8.
Throughout, ``recall'' is reported against its car-only, at-least-ten-surviving-points
denominator, and the two distinct ten-point rules this involves are disambiguated
in Section~4.1.

This scope is narrower than the original proposal in three specific ways, each the
result of an experiment rather than a change of ambition.

\emph{From multi-class to car-only.} The proposal described a pipeline supporting
car, bus, and motorcycle. When the four-class classifier was run on real data it
produced false motorcycle detections --- sequence~08 contains no motorcycles at
all --- and its car recall was poor, and it was replaced by the binary car /
not-car classifier used throughout (Finding~\#20).
The rare classes had too few real examples to learn from and were actively harming
recall; narrowing to cars removed that noise and matches the class distribution of
the data.

\emph{From a CD+EMD completion loss to CD only.} The proposal planned to combine
Chamfer distance with an Earth Mover's distance term to counter the point-clustering
artefacts that Chamfer-only training can produce. That artefact was never the
operative failure mode: the completion outputs were failing as diffuse blobs on real
input because of a domain gap in the input \emph{partiality pattern} --- synthetic
pinhole renders differ from real scan-line LiDAR --- not because of point
non-uniformity (Findings~\#15--\#19). Once that gap was closed with a KITTI-like
partial generator and an inference-normalisation fix (Finding~\#26), Chamfer-only
training was sufficient, and the EMD term (documented but never implemented) was not
needed. The decision here rests on the domain-gap diagnosis, not on the classifier
change above.

\emph{The two-stage-training contribution is withdrawn and reframed.} The proposal
listed a two-stage synthetic-pretraining-then-fine-tuning strategy as a contribution
that would ``dramatically reduce annotation requirements.'' The evidence does not
support that claim: a from-scratch classifier matches the pretrained one, and the
sim-to-real gap is total in both directions, so the synthetic stage is redundant
once real clusters are plentiful rather than annotation-saving (Findings~\#25/\#30).
This is reported as the negative finding it is (contribution~3 above), and the
two-stage training survives only as an ablation.

One point that is \emph{not} a departure from the proposal is worth flagging so it is
not read as one: the proposal was silent on runtime, so the offline framing above is
a scope statement, not a walked-back promise.

\subsection{Thesis outline}
\label{sec:intro-outline}

The remainder of the thesis is organised as follows. Chapter~2 (Background and
Related Work) places the pipeline against the deep-learning segmentation and
point-completion literature and carries the positioning table that defends the
comparison. Chapter~3 (The Detection Pipeline) describes every stage and the
classifier's training. Chapter~4 (Detection Evaluation) sets out the evaluation
protocol and reports the seq-08 results, the ablations, the distance breakdown, and
the runtime. Chapter~5 (The Recall Bottleneck) root-causes the recall shortfall to
clustering fragmentation and reports the chain of negative repair strategies.
Chapter~6 (Point-Cloud Completion) covers the completion method and the debugging
story behind it. Chapter~7 (Completion Evaluation) introduces the donor-frame
metric, shows why pseudo-ground-truth Chamfer is invalid, and reports the
real-data coverage results, the amodal-box utility, the moving-car case, and the
pre-registered held-out replication. Chapter~8 (Discussion and Limitations) names
every weakness the work is subject to, and Chapter~9 (Conclusion and Future Work)
answers the three research questions and lays out what remains.

\bibliographystyle{IEEEtran}
\bibliography{../../references}

\end{document}
